%%%%%%%%%%%%%%%%%%%%%%%%%%%%%%%%%%%%%%%%%%%%%%%%%%%%%%%%%%%%%%%%%%%%%%%%%%% 
% 
% Generic template for TFC/TFM/TFG/Tesis
% 
% By:
%  + Javier Macías-Guarasa.
%    Departamento de Electrónica
%    Universidad de Alcalá
%  + Roberto Barra-Chicote.
%    Departamento de Ingeniería Electrónica
%    Universidad Politécnica de Madrid
% 
% By: + Javier Macías-Guarasa. Departamento de Electrónica Universidad de Alcalá + Roberto Barra-Chicote. Departamento de Ingeniería Electrónica Universidad Politécnica de Madrid
% 
% Based on original sources by Roberto Barra, Manuel Ocaña, Jesús Nuevo, Pedro Revenga, Fernando Herránz and Noelia Hernández. Thanks a lot to all of them, and to the many anonymous contributors found (thanks to google) that provided help in setting all this up.
% 
% See also the additionalContributors.txt file to check the name of additional contributors to this work.
% 
% If you think you can add pieces of relevant/useful examples, improvements, please contact us at (macias@depeca.uah.es)
% 
% You can freely use this template and please contribute with comments or suggestions!!!
% 
%%%%%%%%%%%%%%%%%%%%%%%%%%%%%%%%%%%%%%%%%%%%%%%%%%%%%%%%%%%%%%%%%%%%%%%%%%% 

\chapter{Introducción}
\label{cha:introduccion}

Detrás de cada búsqueda en Internet, cada recomendación en una plantaforma, o cada sistema de inteligencia artificial, existe un 
algoritmo tomando decisiones. Sin embargo, este término suele resultar más difícil de explicar de lo que realmente es, ya que su
concepto se basa en una idea muy sencilla. Un algoritmo es, en esencia, una secuencia de pasos para resolver un problema o 
realizar una tarea de forma ordenada \cite{acevedo}. La historia del pensamiento algorítmico se remonta a la Antigüedad.
Civilizaciones como la antigua Babilonia ya empleaban reglas condicionales para resolver cuestiones legales\footnote{El conocido 
\textit{Código de Hammurabi} \cite{hammurabi} es un ejemplo de ello.}. A lo largo de la historia, este tipo de procedimientos 
también se han utilizado para decidir tratamientos médicos, preparar alimentos y resolver problemas matemáticos, entre muchas 
otras aplicaciones \cite{chabert2012}. No obstante, el término \textit{algoritmo} no apareció hasta el siglo XII, cuando el 
nombre del matemático, astrónomo y geógrafo persa \textit{al-Khwarizmi} fue latinizado como \textit{Algoritmi}, dando origen al
término actual \cite{alkhwarizmi2017}.

Sin embargo, no fue hasta el siglo XIX cuando se desarrolló el primer algoritmo aplicado al mundo de la informática y los 
computadores, con las aportaciones de Ada Lovelace. En 1843, durante la traducción de un artículo sobre la Máquina Analítica
de Charles Babbage \cite{rojas2022}, añadió una serie de notas en las que describió un algoritmo para calcular los números de
Bernoulli \cite{weisstein}, considerado el primer programa diseñado para ser ejecutado por una máquina. Más allá de este logro,
su principal contribución fue anticipar que estas máquinas podrían manipular no solo números, sino cualquier tipo de información
representable mediante símbolos, una idea que sentó las bases de la computación moderna \cite{dasilva2019}.

Las ideas propuestas por Lovelace permanecieron durante décadas como un planteamiento esencialmente teórico. No fue hasta 1936 
cuando Alan Turing introdujo la denominada \textit{máquina de Turing}: un modelo matemático capaz de representar cualquier 
algoritmo mediante un conjunto de reglas simples, ejecutadas de forma secuencial \cite{demol2018}. La máquina de Turing
funcionaba de la siguiente manera: se planteaba como una cinta unidimensional de longitud infinita y dividida en cuadrados, los
cuales contenían cada uno un determinado símbolo, o un vacío. Un cabezal de lectura y escritura recorría la cinta siguiendo un 
conjunto finito de reglas que, en cada ``instrucción'' o paso de control, determinaban las siguientes modificaciones: cambio de
estado, reemplazo del símbolo y desplazamiento a uno de los cuadros adyacentes \cite{rico2021}. Este modelo teórico permitió
definir formalmente el significado de que un problema sea computable, y sentó las bases para estudiar los límites de la 
computación de forma independiente del \textit{hardware}.

Una vez establecido un modelo formal para describir los algoritmos, surgió una nueva cuestión acerca de la eficiencia de cada uno 
de ellos. Aunque cualquier problema computable admite, en principio, una solución algorítmica, no todos los algoritmos presentan 
el mismo rendimiento. Mientras unos resuelven un problema en tiempos razonables, otros requieren un número de operaciones que 
crece de una forma tan rápida que su ejecución resulta impracticable, incluso con la tecnología actual. Esta necesidad de medir 
y comparar la eficiencia de los algoritmos dio lugar al desarrollo de la Teoría de la Complejidad Computacional. Esta teoría 
estudia la cantidad de recursos necesarios para resolver un algoritmo, siendo estos recursos el tiempo\footnote{Número de 
instrucciones ejecutadas.} y el espacio\footnote{Cantidad de memoria empleada.} \cite{cortez2004}.

En la práctica, la complejidad temporal suele recibir más interés que la espacial, ya que, en la mayoría de aplicaciones, el 
tiempo de ejecución destaca como el principal factor limitante del rendimiento. Por tanto, el análisis de algoritmos se centra,
generalmente, en calcular cómo crece el número de instrucciones necesarias a medida que aumenta el tamaño del problema. Para 
cuantificar dicha complejidad temporal se utiliza la notación asintótica, haciendo uso del término \textbf{O grande de}. Esta
notación expresa el orden de crecimiento del tiempo de ejecución de un algoritmo en función del tamaño de entrada, sin tener en 
cuenta constantes o términos de menor orden \cite{phalke2022}.

\begin{table}[ht]
    \renewcommand{\arraystretch}{1.3}
    \caption{Complejidades temporales más comunes.}
    \label{tab:complejidades}
    \begin{center}
        \begin{tabular}{|c|c|}
            \hline
            \textbf{Notación} & \textbf{Nombre} \\
            \hline
            $O(1)$ & Constante \\
            \hline
            $O(\log n)$ & Logarítmica \\
            \hline
            $O(n)$ & Lineal \\
            \hline
            $O(n \log n)$ & Lineal-logarítmica \\
            \hline
            $O(n^k)$ & Polinómica \\
            \hline
            $O(k^n)$ & Exponencial \\
            \hline
            $O(n!)$ & Factorial \\
            \hline
        \end{tabular}
    \end{center}
\end{table}

A medida que fueron desarrollándose algoritmos para resolver problemas cada vez más complejos, surgió la necesidad de 
clasificarlos según su dificultad computacional. Esta clasificación dio lugar a distintas clases de complejidad, las 
cuales agrupan los problemas en función de los recursos necesarios para resolverlos y verificar sus soluciones. La primera 
de ellas es la clase P\footnote{\textit{Polynomial time}.}, y abarca los problemas que se pueden resolver aplicando un
algoritmo en tiempo polinómico, esto es, desde $O(1)$ hasta $O(n^k)$\footnote{Siendo \textit{k} un entero.}. La siguiente 
clase es la NP\footnote{\textit{Nondeterministic Polynomial time}.}, cuyos problemas son aquellos que se pueden verificar\footnote{
Un problema verificable es aquel que, dada una solución candidata, se puede verificar dicha solución en tiempo 
polinómico.} en tiempo polinómico, pero encontrar la solución correcta u óptima requiere un tiempo superior al polinómico, 
aunque inferior al exponencial \cite{cormen2009}. La resolución de estos problemas suele basarse en un enfoque híbrido entre
un algoritmo y una heurística. Finalmente, la clase NP-completo incluye aquellos problemas con una complejidad temporal 
superior a la exponencial. Además, no cuentan con ningún algoritmo eficiente que corra en tiempo polinómico. 

El concepto de NP-completitud apareció por primera vez en 1971, introducido por el científico norteamericano Stephen Cook. En su 
artículo \cite{cook1971} establece que el problema de satisfacibilidad booleana es NP-completo. Este problema consiste en 
determinar si existe alguna asignación de valores de verdad\footnote{Verdadero o falso.} a las variables de una expresión 
lógica booleana, que haga que dicha expresión sea verdadera o satisfacible \cite{academialab}. Un año más tarde, Richard 
Karp, otro científico norteamericano, demostró en su artículo \cite{karp1972} la NP-completitud de otros 21 problemas, entre 
ellos: la cobertura mínima de vértices \cite{vertices}, la coloración de grafos \cite{khanna2000}, el problema de la mochila
\cite{mahapatra2011}, y el problema del camino hamiltoniano \cite{sipser2013}, el cual es un caso especial del problema 
del viajante \cite{johnson2012}.

Es precisamente este último problema clásico, el problema del viajante, el foco de atención de este trabajo. Debido a su alta 
complejidad computacional, la búsqueda de soluciones exactas deja de ser viable cuando aumenta el tamaño del problema. En este 
contexto, toman especial relevancia las heurísticas y metaheurísticas, las cuales son estrategias diseñadas para obtener
soluciones de alta calidad en tiempos de ejecución razonables, pero sin garantizar el óptimo global. 

\section{Objetivos}
\label{sec:obj}

El presente trabajo tiene como objetivo general el estudio, implementación y comparación experimental de cuatro 
metaheurísticas clásicas aplicadas al Problema del Viajante: algoritmos genéticos, optimización mediante colonia de hormigas,
recocido simulado y búsqueda tabú. Para ello, se desarrolla un marco de trabajo propio en Python que permite ejecutar,
evaluar y comparar el rendimiento de estas técnicas sobre un conjunto de instancias estandarizadas del Problema del Viajante,
replicando el diseño experimental de un estudio de referencia en la medida de lo posible.

Con el fin de alcanzar este objetivo general, se plantean los siguientes objetivos específicos:
\begin{itemize}
    \item Estudiar los fundamentos teóricos del Problema del Viajante, su complejidad computacional y sus principales 
    variantes, así como revisar el estado del arte de las técnicas heurísticas y metaheurísticas empleadas para su
    resolución.
    \item Analizar el funcionamiento teórico de las cuatro metaheurísticas seleccionadas, justificando los operadores
    y mecanismos concretos empleados en cada uno de ellas.
    \item Diseñar e implementar, en Python, las cuatro técnicas sobre una representación común del problema, reutilizando
    estructuras y funciones compartidas entre ellas siempre que sea posible.
    \item Seleccionar y justificar los diferentes hiperparámetros de cada algoritmo, apoyándose en la literatura existente
    en lugar de ajustarlos de forma arbitraria, o específicamente para las instancias evaluadas.
    \item Diseñar y ejecutar un experimento comparativo sobre un conjunto de instancias estandarizadas, replicando el diseño 
    de un estudio de referencia y adaptándolo a las limitaciones prácticas de este trabajo.
    \item Analizar y comparar los resultados obtenidos, tanto entre las propias metaheurísticas implementadas como frente a 
    los resultados del estudio de referencia, revisando la calidad de las soluciones, el tiempo de ejecución y la robustez
    de cada técnica.
    \item Extraer conclusiones sobre el comportamiento comparativo de las metaheurísticas estudiadas, y proponer posibles
    líneas de trabajo futuras a partir de los resultados obtenidos.
\end{itemize}


\section{Metodología}
\label{sec:metod}

Para alcanzar los objetivos anteriores, el trabajo se ha estructurado en las siguientes etapas:

\begin{enumerate}
  \item \textbf{Revisión del estado del arte}: se estudian las principales técnicas exactas, heurísticas y metaheurísticas
  empleadas históricamente para resolver el Problema del Viajante, así como estudios comparativos previos entre ellas
  (capítulo~\ref{cha:estudio-teorico}).

  \item \textbf{Estudio del problema}: se formaliza matemáticamente el Problema del Viajante, se analizan sus principales
  variantes y aplicaciones prácticas, y se estudia la representación de una solución candidata sobre la que operarán las
  metaheurísticas (capítulo~\ref{cha:desarrollo}).

  \item \textbf{Estudio teórico de las metaheurísticas}: se analiza en detalle el funcionamiento de cada una de las cuatro
  técnicas seleccionadas, justificando los operadores concretos empleados en cada caso
  (capítulo~\ref{cha:metah}).

  \item \textbf{Diseño e implementación}: se implementan las cuatro metaheurísticas en Python, junto con las funciones
  comunes reutilizadas entre ellas, y se seleccionan y justifican los hiperparámetros de cada algoritmo
  (capítulo~\ref{cha:design}).

  \item \textbf{Diseño experimental y desarrollo de herramientas de apoyo}: se desarrollan los \textit{scripts} necesarios
  para ejecutar el experimento de forma automatizada, recoger los resultados en ficheros auxiliares, calcular las métricas 
  de interés y generar las gráficas comparativas.

  \item \textbf{Ejecución del experimento y análisis de resultados}: se ejecutan las cuatro metaheurísticas sobre el
  conjunto de instancias, y se analizan los resultados obtenidos, comparándolos entre sí y frente al estudio
  de referencia (capítulo~\ref{cha:resultados}).

  \item \textbf{Conclusiones}: se extraen las conclusiones del trabajo realizado, y se proponen posibles líneas de
  continuación a partir de los resultados obtenidos (capítulo~\ref{cha:conc}).
\end{enumerate}

%%% Local Variables:
%%% TeX-master: "../book"
%%% End:


